
\chapter{Процессинговый центр}\label{ch:processing-center}
Каждое архитектурное решение процессингового центра проверяется
двумя вопросами: что произойдёт с задержкой и доступностью
и кто контролирует ключевой материал.

\section{Компоненты процессингового центра}\label{sec:pc-components}

Платёжный шлюз, описанный в предыдущей главе, принимает запрос
ТСП и переводит его в сообщение ISO~8583. Между шлюзом
и карточной сетью стоит процессинговый центр -- система,
которая принимает сообщение, выбирает маршрут, выполняет
криптографические проверки, записывает след операции и позже
превращает авторизацию в клиринговую и расчётную запись.

\definitionnote{Процессинговый центр}{Система, которая авторизует
  карточные операции, формирует клиринг, считает расчётные позиции и
  выполняет криптографические операции. В терминах Закона о
  НПС он часто совмещает роли операционного и платёжного клирингового
центра.}

Процессинг состоит из четырёх основных узлов:
коммутатор авторизации, криптографический контур, клиринговый
модуль и расчётный модуль.

\subsection*{Коммутатор авторизации}

Коммутатор авторизации -- узел, в который
попадает каждое входящее сообщение ISO~8583. Он принимает
запросы от эквайеров или от шлюзов, действующих от их имени,
выбирает маршрут к нужному эмитенту или в карточную сеть
и возвращает ответ. Работает он поштучно и в реальном времени:
допустимая задержка здесь измеряется не секундами,
а десятками миллисекунд.~\cite{visa-core-rules,mc-tpr}

Маршрут внутри коммутатора определяется BIN-таблицей:
для каждого BIN-диапазона известно, через какой канал
(VisaNet, Banknet, ОПКЦ НСПК, прямое подключение к эмитенту)
следует отправить запрос. BIN-таблица обновляется регулярно:
карточные сети публикуют изменения, и процессинг обязан
применять их в установленные сроки, иначе часть транзакций
начнёт уходить не туда и будет отклоняться без вины клиента
или ТСП.~\cite{visa-core-rules,mc-rules}

Коммутатор извлекает
BIN (первые 6--8 цифр PAN) из DE~2, ищет совпадение
в BIN-таблице и определяет карточную сеть (Visa, Mastercard,
Мир). Далее проверяет, принадлежит ли BIN тому же банку,
в котором работает процессинг (on-us). Если on-us --
сообщение маршрутизируется напрямую к авторизационному хосту
эмитента, минуя сеть; interchange при этом не начисляется.
Если off-us -- коммутатор выбирает исходящий канал
(VisaNet, Banknet, ОПКЦ), форматирует сообщение
по Interface Specification конкретной сети (различия
в DE~25, DE~48, subelement layout) и отправляет запрос.
На обратном пути ответ проходит тот же маршрут
в обратном направлении.

Помимо маршрутизации, коммутатор выполняет базовую валидацию
сообщения: проверяет обязательные поля, формат MTI
и корректность bitmap, дополняет запрос данными, которые
требует конкретная сеть, и ведёт журнал транзакций.
Эта запись с точностью до миллисекунд нужна
для отладки и для последующей сверки с клиринговыми
файлами.

Когда внешний маршрут недоступен, включается резервная
авторизация (stand-in processing). Чаще всего stand-in
выполняет сама карточная сеть: VisaNet и Mastercard
предоставляют STIP-сервис, который отвечает за эмитента,
если тот недоступен из-за таймаута или сетевого сбоя.
Процессинговый центр на стороне эмитента тоже может реализовать
локальный stand-in, опираясь на собственные правила: лимит
суммы, BIN-диапазон, историю операций по карте.
В обоих случаях stand-in временно берёт
на себя кредитный риск, который в обычном режиме несёт
эмитент, поэтому его ограничивают малыми суммами
и заранее определёнными категориями транзакций.~\cite{visa-core-rules,mc-rules}

Что происходит, когда эмитент возвращается после сбоя
и обнаруживает, что stand-in одобрил транзакцию, которую
сам эмитент отклонил бы? Ответственность за stand-in-решение
несёт тот, кто его принял. Если stand-in выполняла карточная
сеть (VisaNet STIP, Mastercard STIP), именно сеть берёт
на себя кредитный риск по одобренным операциям
в пределах параметров stand-in: лимита суммы, BIN-списка,
MCC-ограничений. Если stand-in реализован локально
на стороне процессора, риск остаётся на процессоре
или эмитенте -- в зависимости от договора. В обоих случаях
параметры stand-in (макс.\ сумма, разрешённые MCC,
BIN-диапазоны) -- не технический артефакт, а граница
финансовой ответственности: каждый параметр за пределами
нормы -- это потенциальный убыток без покрытия.

\subsection*{Пул HSM}

За криптографию в процессинговом центре отвечают
специализированные аппаратные модули HSM
(Hardware Security Module). Это физически защищённые
устройства, которые хранят ключи в среде, рассчитанной
на сопротивление вскрытию, и выполняют криптографические
операции внутри защищённого периметра
(см. главу~\ref{ch:crypto}).~\cite{pci-pts-hsm,pci-pin-security}

Пул HSM в процессинге решает несколько задач. При PIN-операции
коммутатор извлекает PIN block из DE~52, отправляет его
в HSM вместе с зональным ключом (Zone PIN Key), а HSM
переводит PIN block из формата эквайера в формат эмитента.
При EMV-транзакции тот же контур проверяет криптограмму ARQC:
получает данные операции и ключ эмитента или его производный
вариант, вычисляет ожидаемое значение и сравнивает его
с тем, что пришло от карты. Когда эмитент формирует ответ,
через HSM может быть создана и ARPC
(Authorization Response Cryptogram), которая возвращается
карте для взаимной
аутентификации.~\cite{iso-9564,ansi-x9-24,emvco-book2,emvco-book3}

Пул HSM -- самое жёсткое ограничение для масштабирования
процессинга. Криптографические операции
нельзя разносить по узлам так же свободно, как обычные сервисы
без состояния, а ключевой материал загружается и управляется
по строго регламентированным процедурам. Процессинговые центры строят пул HSM с балансировкой нагрузки
и заранее планируют ёмкость этого контура. Конкретные значения
производительности (операций PIN-перевода в секунду, RSA-подписей
в секунду, ARQC-проверок в секунду) сильно зависят от модели
и редакции HSM, лицензированных функций и параметров ключей;
актуальные числа сверяются по datasheet выбранной модели у её
производителя~\cite{pci-pin-security,pci-pts-hsm,ansi-x9-24}.
Для российского рынка к зарубежным моделям (Thales payShield,
Utimaco, Atalla) добавляется отдельный класс сертифицированных
ФСБ России СКЗИ -- от КриптоПро HSM до решений Аладдин Р.Д.\
и ИнфоТеКС, без которых эксплуатация ГОСТ-контура невозможна.
При росте нагрузки узким местом обычно становится именно HSM,
а не коммутатор или сеть.

\importantnote{%
  HSM остаётся единственным компонентом процессинга, для которого
  горизонтальное масштабирование ограничено физически: каждое
  новое устройство требует церемонии загрузки ключей
  (key ceremony) с принципами разделённого знания
  (split knowledge) и двойного контроля (dual control).
  Планирование
  ёмкости пула HSM должно учитывать текущую нагрузку
  и время, необходимое для ввода в эксплуатацию нового
  устройства.~\cite{pci-pin-security,ansi-x9-24}

}

Управление ключами в пуле HSM устроено иначе, чем в одиночном
устройстве. Мастер-ключ (LMK) нельзя скопировать в открытом виде:
он передаётся через key ceremony, где компоненты вводят несколько
хранителей одновременно. Нарастить ёмкость пула за несколько минут
невозможно -- церемония занимает часы. Синхронизация KCV (Key Check
Value) между устройствами подтверждает, что все HSM пула
работают с одним и тем же ключом, не раскрывая сам ключ. При отказе
одного HSM ключевой материал не теряется: он загружен во все
устройства пула одновременно во время той же церемонии.

\subsection*{Клиринговый модуль}

Клиринговый модуль, в отличие от авторизации, работает
пакетами. Он запускается по расписанию, собирает
авторизованные транзакции за период и добавляет данные,
которых на этапе авторизации ещё не было: финальную сумму
после добавления чаевых (tip adjustment), результат адресной
проверки, параметры DCC-конверсии.

Затем модуль формирует клиринговый файл в формате сети:
TC33 для Visa, IPM для Mastercard
или собственный формат НСПК (см. главу~\ref{ch:reconciliation}).~\cite{visa-core-rules,mc-tpr,nspk-mir-rules}

На этом же этапе система сопоставляет каждую клиринговую
запись с авторизацией по полям: STAN, RRN, сумме
и дате. Так обнаруживаются расхождения: авторизация без
клиринга (orphan authorization), клиринг без авторизации
(late presentment, force post), несовпадение сумм
(частичное подтверждение, tip adjustment). Каждое такое расхождение
означает потенциальную потерю денег, поэтому клиринговый модуль
формирует отчёт об исключениях, который операционная команда
разбирает вручную или по заранее заданным правилам.

\subsection*{Расчётный модуль}

Расчётный модуль считает итоговые нетто-позиции участников:
сколько каждый эмитент должен перечислить и сколько каждый
эквайер должен получить с учётом interchange, комиссий сети
и корректировок вроде чарджбэков, повторных представлений и fees.
Результатом становится набор проводок, который уходит
в расчётную инфраструктуру: для внутрироссийских транзакций
через НСПК и платёжную систему Банка России
(см. главу~\ref{ch:banks}), а для трансграничных сценариев
через расчётные механизмы, предусмотренные правилами
соответствующей
схемы.~\cite{cbr-national-payment-system,cbr-ps,visa-core-rules,mc-switching-services}

Ошибка на расчётном этапе -- прямая финансовая потеря:
именно здесь деньги начинают двигаться между банками.

\section{Производительность и отказоустойчивость}\label{sec:pc-performance}

Требования к задержке и доступности процессингового центра
задаются внешним контрактом сети. Visa публикует maximum time limits
for Authorization Request Response и связывает просрочку
ответа со включением Stand-In Processing; Mastercard
в Transaction Processing Rules задаёт authorization
performance standards для участников сети.~\cite{visa-core-rules,mc-tpr}

\subsection*{Бюджет задержки}

Карточные сети нормируют внешнюю границу ответа, но не дают
универсальной внутренней таблицы \enquote{сколько миллисекунд
должен занимать каждый сервис}. Внутренний бюджет должен
оставлять запас на сетевой путь, обработку у эмитента
и возврат ответа. Конкретная раскладка времени на парсинг,
HSM, логирование и сериализацию зависит от реализации
и измеряется внутри собственного контура. Типичная раскладка
для внутреннего бюджета: парсинг и валидация сообщения
1--3~мс, обращение к HSM (ARQC или PIN-перевод) 3--7~мс,
BIN-lookup и маршрутизация 1~мс, логирование 1--2~мс,
сериализация ответа 1~мс. Итого внутренняя обработка
занимает 7--15~мс; остальное время из бюджета сети
уходит на сетевой путь до эмитента и обратно.

\subsection*{Резервирование и геораспределение}

Процессинговый центр не может позволить себе плановый простой:
даже несколько минут
недоступности -- это тысячи отклонённых транзакций и штрафы
от карточных сетей. Правила сетей не
предписывают единственную допустимую схему, active-active,
active-standby или гибридную, но в любом варианте
оператору нужно переживать отказ площадки без потери
управления трафиком.

Active-active предпочтительнее active-standby по одной причине:
при переключении с passive резерва площадка проходит путь от
холодного состояния до полной нагрузки -- HSM загружает ключи,
BIN-таблицы синхронизируются, очереди восстанавливаются. Каждая
секунда этого пути -- отказанные транзакции. В active-active оба
узла всегда работают под реальной нагрузкой, и при отказе одного
трафик перераспределяется без времени раскрутки. Цена -- риск
split-brain при разрыве репликации. Для авторизации, почти не
имеющей состояния между запросами, последствия кратковременного
split-brain ограничены. Для клиринга и расчёта нужна более жёсткая
гарантия согласованности: авторизационный след можно досверять
по внешним журналам сети, а потерянная клиринговая или расчётная
запись бьёт по денежному движению.

Основная сложность -- синхронизация состояния между ЦОД.
Журнал транзакций,
BIN-таблицы, stand-in лимиты и HSM-ключи должны оставаться
согласованными в обоих центрах.

\importantnote{%
  Контуры авторизации и клиринга/расчёта требуют разных целей
  восстановления: первому важнее быстро восстановиться, второму важнее
  не потерять финансово значимую запись. Поэтому их обычно реплицируют по-разному.

}
\practicenote{%
  Шардинг по BIN-диапазонам остаётся простейшей и наиболее
  распространённой стратегией горизонтального масштабирования
  коммутатора авторизации. BIN однозначно определяет эмитента
  и карточную сеть, а значит, маршрут и набор ключей;
  транзакции с разными BIN не зависят друг от друга
  и могут обрабатываться разными экземплярами коммутатора
  без координации. При необходимости шард можно перенести
  на другой сервер или в другой ЦОД без влияния на остальные
  BIN-диапазоны.

}
\subsection*{Аварийное восстановление}

Георезервирование, то есть размещение ЦОД в разных географических
точках, защищает уже не от отказа отдельного сервера,
а от региональных катастроф: пожара, наводнения, отключения
электричества в районе. Расстояние между центрами всегда
компромиссно. Чем дальше они друг от друга, тем лучше защита
от локальных событий, но тем выше задержка синхронной
репликации. Конкретное расстояние выбирают по сочетанию регионального риска,
качества каналов связи и цены синхронной репликации.

RTO (Recovery Time Objective) -- время восстановления
после полного отказа ЦОД -- для процессинга измеряется
в коротком операционном горизонте: переключение трафика
должно происходить достаточно быстро, чтобы локальный сбой
не превратился в массовый поток отказов. Конкретное значение
определяется архитектурой и обязательствами оператора.

\section{Подключение к платёжным системам}\label{sec:pc-connectivity}

\subsection*{VisaNet}

Подключение к VisaNet, глобальной сети Visa,
определяется правилами Visa и региональными
эксплуатационными регламентами. В публичных документах
Visa отдельно описаны использование VisaNet, требования
к сетевым сертификатам и правила обработки транзакций;
архитектурно VisaNet разделяет BASE~I (авторизация)
и BASE~II (пакетный клиринг).~\cite{visa-core-rules}

Подключение включает канал связи и сертификацию под конкретный
интерфейс сети. Публичные правила Visa оговаривают, что часть деталей,
связанных с безопасностью сети и проприетарной внутренней
обработкой,
в публичной редакции не раскрывается.~\cite{visa-core-rules}

\subsection*{Banknet}

Mastercard использует собственную сеть Banknet.
Публичные материалы Mastercard описывают
коммутационный слой как отдельный контур авторизации,
клиринга и расчёта, а Transaction Processing Rules
одновременно оперируют Authorization Request/0100,
Financial Transaction Request/0200, First Presentment/1240
и форматами IPM.~\cite{mc-switching-services,mc-tpr}

В Mastercard-среде
сосуществуют single- и dual-message сценарии, а клиринговый
слой опирается на IPM и GCMS.~\cite{mc-switching-services,mc-tpr}

\subsection*{ОПКЦ НСПК}

Все внутрироссийские карточные транзакции, независимо
от того, выпущена ли карта внутри российской системы Мир
или относится к международной схеме, работающей во внутреннем
контуре через российскую инфраструктуру, упираются
в НСПК. Банк России прямо указывает, что внутрироссийская
обработка операций международных платёжных систем ведётся
через НСПК; сама НСПК публично описывает себя как оператора
ПС Мир и ОПКЦ НСПК.~\cite{cbr-national-payment-system,nspk-about-company}

Для российского банка выход во внутренний
карточный контур означает подключение
к ОПКЦ НСПК и соблюдение правил ПС Мир и связанных
технических документов НСПК.~\cite{nspk-mir-rules}

\subsection*{СБП}

Подключение к СБП архитектурно отличается от карточного
процессинга. СБП работает как сервис платёжной системы Банка России,
для которого Банк России публикует отдельный порядок
подключения кредитных организаций, а НСПК публикует
материалы по сценариям для бизнеса, оплате по QR, кнопке, NFC
и возвратам для бизнеса.~\cite{cbr-sbp-actions,sbp-main,sbp-faq-business}

Процессинг, который поддерживает и карты, и СБП,
совмещает в одной инфраструктуре два разных класса потоков:
карточный контур на сообщениях ISO~8583 и сервис перевода средств
по правилам СБП.~\cite{cbr-sbp-actions,sbp-main}

\subsection*{Перегрузка и управляемая деградация}

Когда карточная сеть не отвечает, процессинговый центр не может
просто ждать: тысячи новых авторизаций продолжают поступать, и
очередь быстро переполнится.

Типовая защита -- таймауты внешних вызовов и автоматическое
размыкание канала при массовых сбоях. Точные пороги
задаются конкретной реализацией процессинга.

\section{Лицензирование и сертификация}\label{sec:pc-licensing}

Процессинговый центр не может работать без лицензии
регулятора, сертификации у карточных сетей и аудита PCI DSS.

\subsection*{Требования Банка России}

Закон о НПС разводит роли операционного центра (ОЦ),
платёжного клирингового центра (ПКЦ), расчётного центра (РЦ)
и других операторов услуг платёжной инфраструктуры (ОУПИ);
именно в этой терминологии процессинговый контур получает
свой юридический статус, а сами ОУПИ включаются в реестр
операторов услуг платёжной инфраструктуры, который ведёт
Банк России.~\cite{fz-161,cbr-ps-glossary}\footnote{Точные
реквизиты порядка ведения реестра ОУПИ и состав публикуемых
сведений сверяются по актуальной редакции на сайте Банка
России перед публикацией главы.}

Актуальные требования к защите информации задаются на нескольких
уровнях. Защиту информации при осуществлении переводов денежных
средств регулирует Положение Банка России от 17.08.2023 № 821-П,
вступившее в силу 1 апреля 2024 года~\cite{cbr-821p}; оно пришло
на смену Положению № 719-П (действовало с 2022 года), которое,
в свою очередь, заменило ранее применявшееся Положение № 382-П
от 09.06.2012.\footnote{Ссылки на 719-П и 382-П в действующих
документах процессинга и в типовых договорах с участниками
платёжных систем подлежат миграции на 821-П.}
Положение Банка России от 30.01.2025 № 851-П, действующее
с 29 марта 2025 года, закрепляет обязанности банков по защите
информации в целях противодействия переводам без согласия
клиента и заменило ранее применявшееся
Положение № 683-П~\cite{cbr-851p}.
К защите информации в платёжной системе Банка России и в системе
передачи финансовых сообщений (СПФС) применяется отдельная
норма -- Положение Банка России от 25.07.2022
№ 802-П~\cite{cbr-802p}.

Технической базой для всех перечисленных требований служит
национальный стандарт ГОСТ Р 57580.1-2017 (требования к защите
информации) и связанный с ним ГОСТ Р 57580.2-2018 (методика
оценки соответствия), которые задают три уровня защиты --
минимальный, стандартный и усиленный -- и соответствующий
набор обязательных мер~\cite{gost-57580-1,gost-57580-2}.
К процессинговому центру, обрабатывающему персональные данные
российских граждан, дополнительно применяется требование
152-ФЗ о размещении и первичной обработке таких данных на
территории Российской Федерации (часть~5 статьи~18); с 1~июля
2025~года действует обновлённая редакция статьи, введённая
Федеральным законом от 28~февраля 2025~года № 23-ФЗ.~\cite{fz-152}

Для участия в национальной платёжной системе \enquote{Мир}
процессинговый контур проходит сертификацию ПЦМ (процессинговый
центр Мир) -- программу сертификации процессинга для ПС~\enquote{Мир},
которую проводит НСПК как оператор ПС~\cite{nspk-mir-rules};
для участия в СБП -- интеграцию с ОПКЦ СБП на стороне
НСПК~\cite{cbr-sbp-rules}. Эти контуры дополняют, а не
заменяют требования Банка России.

\subsection*{Сертификация у карточных сетей}

Каждая карточная сеть проводит собственную программу
сертификации для процессинговых центров, подключающихся
к её инфраструктуре. В публичных правилах Visa и Mastercard
зафиксированы сами рамки участия в сети: правила использования
VisaNet, требования к обработке транзакций, составу данных,
клирингу и расчёту.~\cite{visa-core-rules,mc-rules,mc-tpr}

Сертификация охватывает корректность форматов и обязательных
элементов данных, криптографические процедуры и устойчивость
операционного контура к сбоям. Ни Visa, ни Mastercard в открытых
редакциях не раскрывают полную матрицу тестов, связанных
с безопасностью. Детальные приёмочные сценарии живут
в закрытых руководствах и становятся доступны участнику уже
на этапе подключения.~\cite{visa-core-rules,mc-rules}

\subsection*{PCI DSS}

Процессинговый центр почти неизбежно попадает в CDE
(cardholder data environment): PCI SSC отдельно разводит
cardholder data и sensitive authentication data, а для PIN,
PIN block и связанных процедур управления ключами действует
самостоятельный стандарт PCI PIN
Security.~\cite{pci-scoping-guidance,pci-faq-1335,pci-pin-security}

Требования к подтверждению соответствия зависят
от роли и категории сервис-провайдера в программе конкретного
бренда, а не только от формальной отнесённости к Level~1. У Visa для Level~1 service providers и VisaNet
processors используется ROC by QSA, а Mastercard для
регистрируемых сервис-провайдеров требует ежегодного
подтверждения соответствия PCI DSS и для применимых
категорий ежеквартальных
ASV-сканов.~\cite{visa-pci-validation-best-practice,mc-service-provider-categories-pci,pci-faq-1234-asv}
